\section{Conclusion}

Context layers work. The open question was which part of one does the work, and
it matters because the parts cost different amounts: semantic prose is written
once per domain and often already exists, retrieval scaffolding is engineering
paid once, and pre-computed macros are paid per metric forever.

Holding the tool surface, the retrieval instruction, the table allow-list and
the operation rules byte-for-byte fixed and emptying only the prose,
\textbf{the content step dominates the scaffolding step on every model} --- by
$1.8{\times}$ on the strongest model, $6.6{\times}$ on the weakest, and
entirely on the one where scaffolding does nothing at all.
Delivery matters as much as possession: the same governing rule sat in both
arms, and only the contract arm applied it --- visibly, in the SQL it
submitted, where it writes all six load-bearing clauses up to 98\% of the time
against the prompt arm's 4\%.

Because the contract's declared expressions compile to a layer that reproduces
gold exactly, each arm's accuracy can be read as efficiency against a proven
ceiling. Read that way the contract arm recovers 60.8\%, 62.7\% and 94.9\% of
what it was handed: \textbf{the cost of requiring an agent to derive, rather
than handing it a pre-computed macro, is not structural but transient}, and
already small at the frontier.

We also report two corrections. On two models the scaffolding step was
indistinguishable from zero and the contract's margin over a prompt appeared to
shrink with capability; we stated both, and the third model overturned both.
\textbf{A null across two adjacent models is weak evidence of absence}, which
is the methodological result of this paper and the one we would most want a
reader to carry away.

The same transcripts carry a second, narrower result. Declared rules that the
model merely reads stopped every mutating statement it would otherwise have
submitted --- 144 to zero, on the two models that reached for the idiom at all
--- while the enforcement layer standing behind those rules never had to fire.
On the third model no arm attempted a write, so the hazard disappeared before
the mechanism could be tested on it.

Three limits bound all of this. The gain appears only inside the domain the
contract describes. The benchmark is a single one, and a concentrated one. And
no run
repeats, so nothing here estimates within-condition variance.

\paragraph{Future work.} Three things, in order of what they would settle.
Repeat runs would put a variance estimate under every number here; that is the
largest gap, and largest for the frontier run, whose temperature could not be
pinned. A second frontier model would test every capability claim in this
paper, all of which now rest on one, which is the exact condition under which a
claim here has already broken once. And an executable metric layer, built
\emph{after} this measurement instead of before it, would separate the
declarative contribution from the executable one the way this paper separates
content from scaffolding.

What we would defend today is narrower than what two models supported: the
content in a context layer does most of the work, the scaffolding does some,
and the cost of making an agent derive a query instead of looking it up is
falling fast enough that it should not decide the design.

\section*{Artifact Availability}

The harness, the frozen contract and its mechanically derived hollow variant,
all 4{,}812 transcripts including model reasoning, and the raw result rows are
available at
\url{https://github.com/flyersworder/agentic-data-contracts}. Every row records
the commit, contract digest, gold-set hash, scorer identity, provider endpoint,
quantization, token counts, cost, and a post-run database integrity check.
